% 本文件由 LaTeX 排版 Agent 根据有效 translation.json 自动生成。
% author=Apocrypha; work_id=1007; title=《亚巴郎遗嘱》

\chapter{《亚巴郎遗嘱》}

% structure: p0001 source=h1 target=omitted reason=page-title-already-rendered-by-parent

% structure: p0002 source=h2 target=section reason=relative-heading-rank; base=h2

\section{版本一}

1. 亚巴郎度了他生命的年岁，共九百九十五岁。他一生以恬静、温和和正义度过了所有年岁，这位义人极其好客；因为他在玛默勒橡树处的交叉路口支搭帐幕，接待每一个人，无论贫富、君王和统治者、残疾和无助者、朋友和陌生人、邻人和旅客，虔诚、至圣、正义而好客的亚巴郎都一样款待。然而，连他也逃脱不了那普遍、无情而痛苦的死亡命运，以及人生无常的终结。因此，上主天主召唤他的总领天使弥额尔，对他说：\Quote{下去，总领天使弥额尔，到亚巴郎那里，向他谈论他的死亡，好叫他安排他的事务，因为我已祝福他如同天上的星辰和海边的沙粒，他享有长寿和许多财产，且变得极其富有。此外，他在一切善事上比众人更正义，好客而仁爱，直到生命尽头；但你去吧，总领天使弥额尔，到我亲爱的朋友亚巴郎那里，向他宣布他的死亡，并这样向他保证：你此时要离开这虚幻的世界，舍弃肉身，到你的主那里，与善人同在。}\par

2. 总领天使从天主面前出发，下到亚巴郎那里，到玛默勒橡树处，在附近的田里遇见义人亚巴郎，他正坐在牛轭旁耕田，与马塞克的儿子们和其他仆人一起，共十二人。看，总领天使来到他那里，亚巴郎看见总领天使弥额尔从远处走来，酷似一位极其英俊的武士，就起身照常迎接他，接待并款待所有陌生人。总领天使向他请安说：\Quote{万福，最可敬的父亲，天主所拣选的义人灵魂，天上者的真子。}亚巴郎对总领天使说：\Quote{万福，最可敬的武士，明亮如太阳，在众人子中最美丽；欢迎你，因此我恳求你，告诉我你青春的年岁来自何处；教导我这恳求者，你的美丽来自哪支军队、哪段旅程。}总领天使说：\Quote{我，义人亚巴郎啊，来自那座大城。我被大王派遣来取代他的一位好友的位置，因为王召见了他。}亚巴郎说：\Quote{请来，我主，与我一同到我的田里。}总领天使说：\Quote{我来。}于是他们进入耕田，在人群旁边坐下。亚巴郎对他的仆人，马塞克的儿子们说：\Quote{到马群那里，牵两匹安静、温顺而驯良的马，好让我和这位客人骑上。}但总领天使说：\Quote{不，我主亚巴郎，不要让他们牵马来，因为我从不骑任何四足牲畜。难道我的王不是富有大量货物，拥有统治人和各类牲畜的权力吗？但我从不骑任何四足牲畜。那么，我们走吧，义人的灵魂，轻步直到我们到达你的家。}亚巴郎说：\Quote{阿们，但愿如此。}\par

3. 当他们从田里往他家走去时，路旁有一棵柏树，因主的命令，那树用人的声音呼喊说：\BibleQuote{圣！圣！圣！上主天主，他呼召那些爱他的人。}但亚巴郎隐瞒了这奥秘，以为总领天使没有听见树的声音。走近屋子后，他们坐在庭院里，依撒格看见天使的面容，对他母亲撒辣说：\Quote{我的母亲夫人，看，与我父亲亚巴郎同坐的那人，不是地上居民的族类之子。}依撒格就跑过去，向他请安，俯伏在\OriginalTerm{无形体者}{Incorporeal}脚前，无形体者祝福他说：\Quote{上主天主必赐给你他向你父亲亚巴郎和他的后裔所应许的，也必赐给你你父亲和母亲宝贵的祈祷。}亚巴郎对他的儿子依撒格说：\Quote{我儿依撒格，从井里打水，用器皿盛来给我，好给这位客人洗脚，因为他累了，从远方来到我们这里。}依撒格就跑到井边，用器皿打了水，拿来给他们；亚巴郎上前为总领天使弥额尔洗脚，亚巴郎的心被感动，为客人哭泣。依撒格见父亲哭泣，也哭了；总领天使见他们哭泣，也同他们一起哭。总领天使的泪落在器皿里盆中的水上，变成了宝石。亚巴郎见了这奇事，惊讶不已，就悄悄地拿起宝石，隐瞒了这奥秘，把它藏在心中。\par

4. \Person{亚巴郎}对他的儿子\Person{依撒格}说：\Quote{我亲爱的儿子，去屋里的内室并装饰它。为我们铺两张床，一张为我，一张为今天与我们同住的这位客人。为我们预备座位、灯台和一张桌子，摆满各种佳肴。装饰内室，我儿，并在我们下面铺上细麻、紫色和上等细麻。在那里焚烧各种珍贵和上好的香料，从花园里带来芬芳的植物，使我们的屋子充满香气。点燃七盏装满油的灯，好让我们欢欣，因为今天与我们同住的这位客人比君王或统治者更荣耀，他的容貌超越所有世人。} \Person{依撒格}便妥善预备了一切，\Person{亚巴郎}带着总领天使\Person{米迦勒}进入内室，两人都坐在床上，在他们中间放了一张桌子，摆满各种佳肴。随后，那总领天使起身出去，仿佛因肚腹的约束要去排泄，眨眼之间升到天上，站在主面前，对他说：\Quote{主上，愿你的权能知道，我无法提醒那义人关于他的死亡，因为我在地上没有见过像他这样的人，富有怜悯、好客、公义、诚实、虔诚、远离一切恶行。现在，主，你知道我无法提醒他关于他的死亡。} 主说：\Quote{下去，总领天使\Person{米迦勒}，到我的朋友\Person{亚巴郎}那里，凡他对你说的，你都去做，凡他吃的，你与他同吃。我将派遣我的圣神到他儿子\Person{依撒格}心中，将死亡之念放在\Person{依撒格}心里，以致他甚至在梦中能见到他父亲的死亡，\Person{依撒格}将述说那梦，你将解释它，他本人便会知道他的结局。} 那总领天使说：\Quote{主，所有天上的神灵都是无形体的，既不饮也不食，这个人却在我面前摆了一张桌子，摆满了地上的、可朽坏的各种佳肴。现在，主，我该怎么办？我如何回避与他同坐一桌？} 主说：\Quote{下去到他那里，不必为此忧虑，因为当你与他同坐时，我将派遣一个吞噬的灵到你身上，它必从你手中并经由你的口吞尽桌上的一切。你要同他一起在所有事上欢喜，只是你要妥善解释那异象中的事，好让\Person{亚巴郎}知道死亡的镰刀和生命的不确定结局，并妥善处置他的一切财物，因为我降福他超过海沙和天上的星辰。}\par

5. 然后，那总领天使下到\Person{亚巴郎}的家中，与他一同坐在桌前，\Person{依撒格}伺候他们。用毕晚餐，\Person{亚巴郎}照常祈祷，那总领天使也同他一起祈祷，然后各人在自己的床上躺下睡觉。\Person{依撒格}对他的父亲说：\Quote{父亲，我也愿同您一起在这屋里睡觉，好让我也能聆听您的谈论，因为我喜爱听这位有德之人谈话的高妙之处。} \Person{亚巴郎}说：\Quote{不，我儿，你回自己的屋里，在自己的床上睡觉，免得我们打扰这个人。} 于是\Person{依撒格}从他们那里领受了祝福，并祝福了他们，然后回自己的屋里，在自己的床上躺下。但主将死亡之念如同梦境一般投入\Person{依撒格}心中，约夜间的第三时辰，\Person{依撒格}醒来，从床上起来，跑到他父亲与总领天使一同睡觉的屋里。\Person{依撒格}一到门口便哭喊说：\Quote{我的父亲\Person{亚巴郎}，起来，快给我开门，让我进来，挂在您的颈上，拥抱您，免得他们将您从我这里带走。} \Person{亚巴郎}于是起来给他开门，\Person{依撒格}进来挂在他的颈上，放声大哭。\Person{亚巴郎}因此内心激动，也放声大哭，那总领天使见他们哭泣，也哭泣起来。\Person{撒辣}在自己的房间听到他们的哭声，跑来找他们，见他们拥抱哭泣。\Person{撒辣}哭着说：\Quote{我的主\Person{亚巴郎}，您为何哭泣？请告诉我，我的主，今天与我们同住的这位兄弟是否带来了您的侄儿\Person{罗特}的死讯？您是为这事如此悲伤吗？} 那总领天使回答她说：\Quote{不，我的姊妹\Person{撒辣}，并非如你所说，而是你的儿子\Person{依撒格}，我想是做了一个梦，来到我们这里哭泣，我们见他这样，心中感动，也哭了。}\par

6. 那时，\Person{撒辣}听见那总领天使谈话的高妙之处，立刻知道这是主的天使在说话。\Person{撒辣}便示意\Person{亚巴郎}到门边，对他说：\Quote{我的主\Person{亚巴郎}，您知道这人是谁吗？} \Person{亚巴郎}说：\Quote{我不知道。} \Person{撒辣}说：\Quote{我的主，您还记得那天从天上来的三位男子，他们在\Place{玛默勒}橡树旁的帐篷里受到我们的款待，那时您宰了一只无瑕疵的山羊羔，在他们面前摆了一张桌子。那些肉被吃掉后，那山羊羔又活过来，欢乐地吸吮母亲的乳。我的主\Person{亚巴郎}，您难道不知道，他们曾应许给我们\Person{依撒格}作为腹中的果实？这三位圣者中的一位正是他。} \Person{亚巴郎}说：\Quote{哦，\Person{撒辣}，你这话说得对。愿赞美和颂赞归于我们的天主和圣父。因为傍晚我在盆里洗他的脚时，心里想，这正是我那时所洗的三位中的一人的脚；那时落在我盆里的他的泪水化成了宝石。} 他从怀中取出宝石，交给\Person{撒辣}，说：\Quote{你若不信我，现在看看这个。} \Person{撒辣}接过来，俯伏致敬说：\Quote{愿光荣归于天主，他向我们显明了奇妙的事。} \Person{撒辣}又说：\Quote{现在，我的主\Person{亚巴郎}，要知道我们中间正有一件启示发生，无论是祸是福！}\par

7. \Person{亚巴郎}离开\Person{撒辣}，进入内室，对\Person{依撒格}说：\Quote{我亲爱的儿子，来，告诉我实情：你看见了什么？遭遇了什么事，为何如此匆忙地来到我们这里？}\Person{依撒格}回答说：\Quote{我主，今夜我看见太阳和月亮在我头顶上，用它的光芒环绕着我，给我光明。当我看着并欢喜时，我看见天开了，一个发光的人从天而降，比七个太阳还要明亮。这个像太阳的人来了，从我头上取走了太阳，然后升到天上他来的地方。但我非常悲伤，因为他取走了我的太阳。过了一会儿，当我仍忧伤痛苦时，我看见这个人第二次从天上出来，又从我头上取走了月亮。我大哭，向那发光的人呼求说：\InnerQuote{我主，请不要夺走我的荣耀；可怜我，听我说：如果你从我这里取走太阳，就把月亮留给我。}他说：\InnerQuote{任凭它们被带到上面的君王那里，因为他希望它们在那里。}他便把它们从我这里取走了，却把光芒留给了我。}天军之帅说：\Quote{正义的\Person{亚巴郎}，请听：你儿子看见的太阳便是你，他的父亲；月亮则是他的母亲\Person{撒辣}。那从天上降下的发光之人，便是从天主派来取你正义灵魂的那一位。现在，最受尊崇的\Person{亚巴郎}，你要知道，此时你将离开这尘世的生命，前往天主那里。}\Person{亚巴郎}对天军之帅说：\Quote{真是最奇妙的奇事！现在，你就是那位要从我身上取走我灵魂的那位吗？}天军之帅对他说：\Quote{我是天军之帅\Person{弥额尔}，侍立在上主面前，我被派到你这里来提醒你关于你的死亡，然后我将按所吩咐的回到他那里。}\Person{亚巴郎}说：\Quote{现在我知道你是上主的天使，被派来取我的灵魂，但我不愿跟你去；你只管照你所受命的去做吧。}\par

8. 天军之帅听了这些话，立即隐没，升到天上，站在天主面前，述说了他在\Person{亚巴郎}家中看见的一切；天军之帅也对上主这样说：\Quote{你的朋友\Person{亚巴郎}这样说：\InnerQuote{我不愿跟你去，你只管照你所受命的去做吧。}现在，全能的上主啊，你的荣耀和不朽的国度有什么命令吗？}天主对天军之帅\Person{弥额尔}说：\Quote{你再到我朋友\Person{亚巴郎}那里去一次，对他说：\InnerQuote{上主你的天主这样说：是我领你进入应许之地，是我祝福你多于海沙和天上星辰，是我开启了\Person{撒辣}不育的子宫，在老年赐给你\Person{依撒格}作为腹中的果实。我实在告诉你：我必赐福给你，使你的后裔繁多，并将你向我祈求的一切赐给你，因为我是上主你的天主，除我以外没有别的神。告诉我，你为什么背叛我？为什么你心中忧愁？你为什么背叛我的总领天使\Person{弥额尔}？难道你不知道所有从\Person{亚当}和\Person{厄娃}而来的人都死了，没有一位先知逃脱了死亡吗？没有一位作王的不朽；你的祖先没有一人逃脱了死亡的奥秘。他们都死了，都进入了阴府，都被死亡的镰刀收聚了。但我没有派死亡到你那里，没有让任何致死的疾病临到你，没有允许死亡的镰刀遇见你，没有让阴府的网罗缠住你，我从未愿意你遭遇任何凶恶。而是为了美好的安慰，我派了我的天军之帅\Person{弥额尔}到你那里，好让你知道你要离开世界，料理你的家和你所有的一切，并祝福你亲爱的儿子\Person{依撒格}。现在你要知道，我做这事并非愿意使你忧伤。那么，你为什么对我的天军之帅说：\InnerQuote{我不愿跟你去}？你为什么这样说？难道你不知道，如果我让死亡来到你身上，那时我就要看看你是否会来吗？}}\par

9. 天军之帅领受了上主的训诲，下到\Person{亚巴郎}那里。义人一看见他，便仆倒在地，如同死人一般。天军之帅将他从至高者那里听到的一切都告诉了他。于是神圣而正义的\Person{亚巴郎}起身，泪流满面，俯伏在那无形体的脚前，恳求他说：\Quote{我恳求你，天上万军之帅，既然你完全屈尊亲自来到我这罪人、你各方面不配的仆人这里，我现今仍然恳求你，天军之帅啊，请再将我的话语带到至高者那里，你要对他说：\InnerQuote{你的仆人\Person{亚巴郎}这样说：主，上主！在我向你祈求的每件工作和话语上，你都垂听了我，并完成了我的全部意愿。现在，上主，我不抗拒你的大能，因为我也知道我不是不朽的，而是必朽的。既然万物都服从你的命令，在你大能的面前畏惧战栗，我也畏惧，但我向你求一件事。现在，主上，求听我的祈祷：当我还在这个身体内时，我渴望看见所有有人居住的大地，以及你一言所建立的一切受造物。我看见了这些，然后若我离开生命，我就没有忧愁了。}}于是天军之帅又回去，站在天主面前，告诉他一切，说：\Quote{你的朋友\Person{亚巴郎}这样说：\InnerQuote{我渴望在我死前在有生之年看见全地。}}至高者听了这话，又命令天军之帅\Person{弥额尔}，对他说：\Quote{取一朵光云，以及那些掌管战车的天使，下去，把正义的\Person{亚巴郎}放在革鲁宾的战车上，将他高举到天上的空中，使他能看见全地。}\par

10. 总领天使弥额尔下来，让亚巴郎乘上革鲁宾的战车，将他高举到天上的空中，带领他连同六十位天使一同在云上，亚巴郎乘着战车越过全地。亚巴郎看见那日世界的样子：有人耕种，有人驾车，一处有人牧放羊群，另一处有人夜间看守，还有跳舞、游戏和弹琴，又一处有人争辩打官司，别处有人哭泣并追念死者。他也看到新婚之人受到尊敬，总之，他看见了世上所行的一切事，无论善恶。亚巴郎经过他们时，看见有人手持刀剑，手中挥舞着锋利的刀剑。亚巴郎问总领天使说：\Quote{这些人是谁？}总领天使说：\Quote{这些是盗贼，他们打算杀人、偷窃、焚烧和毁灭。}亚巴郎说：\Quote{主啊，主啊，请听我的声音，命令野兽从林中出来吞噬他们。}他正说话时，就有野兽从林中出来吞噬了他们。他又在另一处看见一个男人与一个女人彼此行淫，便说：\Quote{主啊，主啊，命令地裂开吞没他们。}地立即裂开吞没了他们。他又在另一处看见有人挖穿房屋，搬走别人的财物，便说：\Quote{主啊，主啊，命令火从天降下烧灭他们。}他正说话时，就有火从天降下烧灭了他们。随即有声音从天上向总领天使这样说：\Quote{总领天使弥额尔，命令战车停下，使亚巴郎转向，以免他看见全地；因为他若看见所有活在邪恶中的人，便会毁灭一切受造物。因为，看，亚巴郎没有犯罪，也不怜悯罪人；但我创造了世界，不愿毁灭其中任何一个，而是等待罪人之死，直到他悔改而生活。但你将亚巴郎带到天堂的第一道门，使他在那里看见审判和报应，并为他所毁灭的罪人灵魂而悔改。}\par

11. 于是弥额尔转动战车，将亚巴郎带到东方，到了天堂的第一道门。亚巴郎看见两条路：一条狭窄而紧缩，另一条宽阔而宽敞；他又看见两道门：一道门宽阔，在宽阔的路上；另一道门狭窄，在狭窄的路上。在那两道门外，他看见一个人坐在镀金的宝座上，那人的面貌可怕，如同主一般。他们看见许多灵魂被天使驱赶，从宽门被引入；另有一些灵魂，为数稀少，由天使通过窄门带走。当那位坐在金宝座上的奇妙者看见少数人从窄门进入，多数人从宽门进入时，他立刻扯自己的头发和胡须两边，从宝座上扑倒在地，哭泣哀号。但当他看见许多灵魂从窄门进入时，便从地上起来，满心欢喜地坐在宝座上，欢欣踊跃。亚巴郎问总领天使说：\Quote{我主总领天使，这位最杰出的人是谁？他装饰着如此荣耀，有时哭泣哀号，有时欢喜踊跃？}那无形者说：\Quote{这是首造者亚当，他处于这样的荣耀中。他观看世界，因为众人都是从他而生。当他看见许多灵魂从窄门进入时，便起来坐在宝座上欢欣踊跃，因喜乐而雀跃，因为这道窄门是义人的门，引人进入生命，从它进入的人进入乐园。因此，首造者亚当欢喜，因为他看见灵魂得救。但当他看见许多灵魂从宽门进入时，便拔自己的头发，扑倒在地，痛哭哀号，因为宽门是罪人的门，引人进入灭亡和永罚。因此，首造者亚当从宝座上跌倒，为罪人的灭亡而哭泣哀号，因为丧失的人很多，得救的人很少，因为在七千人中间几乎找不到一个灵魂得救，是公义无玷的。}\par

12. 他还在对我说这些话时，看，两位天使，面貌如火，心肠无情，目光严厉，驱赶着成千的灵魂，毫不怜悯地用火鞭抽打他们。天使抓住一个灵魂，把所有的灵魂都赶进宽门去受毁灭。我们也随着天使进去，进入那宽门之内。在两门之间，立着一个可怕的宝座，由可怕的、像火一样发光的晶体做成，上面坐着一位杰出的人，明亮如太阳，如同天主子一般。他面前立着一张桌子，像晶体一样，全是金子和细麻，桌上放着一本书，厚度六肘，宽度十肘。左右各站着两位天使，拿着纸、墨和笔。桌子前坐着一位光明天使，手中持着一架天平；他的左边坐着一位全身如火、无情而严厉的天使，手中持着一支号角，号角内装有耗尽的火，用以试验罪人。坐在宝座上的那位杰出者亲自审判和定夺灵魂，左右两位天使记下，右边的记下义行，左边的记下恶行。那在桌子前持天平的天使称量灵魂，那持火的天使用火试验灵魂。亚巴郎问总领天使弥额尔说：\Quote{我们所看见的这些是什么？}总领天使说：\Quote{圣洁的亚巴郎，你所看见的这些是审判和报应。看，那天使手中拿着一个灵魂，带到审判官面前。审判官对侍奉他的一个天使说：\InnerQuote{给我打开这本书，找出这灵魂的罪。}他打开书，发现它的罪和义行相等，便没有把它交给刑罚者，也没有交给得救者，而是放在中间。}\par

13. \Person{亚巴郎}说：\Quote{\Person{我主}·\OriginalTerm{主帅}{chief-captain}，这个最奇妙的审判者是谁？记录的天使是谁？像太阳一样、拿着天平的天使是谁？拿着火的火天使是谁？}\OriginalTerm{主帅}{chief-captain}说：\Quote{至圣的\Person{亚巴郎}，你看见坐在宝座上的可怕的人吗？这是第一个受造的人\Person{亚当}的儿子，名叫\Person{亚伯尔}，被恶者\Person{加音}杀害。他这样坐着审判所有受造之物，分辨义人和罪人。因为\Person{天主}曾说：\InnerQuote{我不审判你们，但每个由人所生的人都要被审判。}因此，他把审判权交给他，审判世界，直到他伟大而光荣的来临。然后，正义的\Person{亚巴郎}，那将是完美而不可更改的审判和报应，永恒而不变，无人能改变。因为每个人都来自第一个受造者，因此他们首先在这里由他的儿子审判，在第二次来临时，他们将由\Person{以色列}十二支派审判，每一个气息和每一个受造物。但第三次，他们将由万有的主\Person{天主}审判，那时，审判的结局近了，判决可怕，无人能解救。现在，通过三个法庭，世界的审判和报应得以完成，因此，一件事不能凭一两个证人最后确定，而应凭三个证人确定一切。右边和左边的两位天使，他们是记录罪恶和义行的，右边的那位记录义行，左边的那位记录罪恶。像太阳一样、手里拿着天平的天使，是总领天使\Person{多基耳}，公义的称量者，他用\Person{天主}的正义来衡量义行和罪恶。拿着火、毫无怜悯的天使，是总领天使\Person{普鲁耳}，他掌管火，用火试验人的行为，如果火烧毁了某人的行为，审判天使立刻抓住他，把他带到罪人的地方，一个极其痛苦的惩罚之地。但如果火认可某人的行为，没有烧毁它，那人就被称为义，正义的天使就带他上去，使他得救，进入义人的份。因此，最正义的\Person{亚巴郎}，所有人在一切事上都要经火和天平试验。}\par

14. \Person{亚巴郎}对\OriginalTerm{主帅}{chief-captain}说：\Quote{\Person{我主}·\OriginalTerm{主帅}{chief-captain}，天使手里拿着的那个灵魂，为什么被判定放在中间呢？}\OriginalTerm{主帅}{chief-captain}说：\Quote{正义的\Person{亚巴郎}，请听。因为审判者发现它的罪恶和义行相等，既没有把它交审判，也没有使它得救，直到万有的审判者来到。}\Person{亚巴郎}对\OriginalTerm{主帅}{chief-captain}说：\Quote{那灵魂要得救，还缺少什么呢？}\OriginalTerm{主帅}{chief-captain}说：\Quote{如果它获得一个超过其罪恶的义行，它就能进入救恩。}\Person{亚巴郎}对\OriginalTerm{主帅}{chief-captain}说：\Quote{\OriginalTerm{主帅}{chief-captain}\Person{弥额尔}，请到这里来，让我们为这个灵魂祈祷，看\Person{天主}是否会听我们。}\OriginalTerm{主帅}{chief-captain}说：\Quote{阿们，但愿如此。}他们为灵魂祈祷和恳求，\Person{天主}听了他们。当他们从祈祷中站起来时，没有看见灵魂站在那里。\Person{亚巴郎}对天使说：\Quote{你放在中间的那个灵魂在哪里？}天使回答说：\Quote{因你正义的祈祷，它得救了。看，一个光明的天使已经把它带走，提到\Place{乐园}里去了。}\Person{亚巴郎}说：\Quote{我光荣至高\Person{天主}的名和他的无量慈悲。}\Person{亚巴郎}对\OriginalTerm{主帅}{chief-captain}说：\Quote{总领天使，我恳求你，俯听我的祈祷，让我们再呼求主，恳求他的怜悯，为那些我以前在愤怒中诅咒和毁灭的灵魂祈求他的仁慈——那些被大地吞没、被野兽撕裂、被火焚烧的灵魂，都是因我的言语。现在我知道，我在主我们的\Person{天主}面前犯了罪。所以，来，天上的军队的\OriginalTerm{主帅}{chief-captain}\Person{弥额尔}，来，让我们含泪呼求\Person{天主}，使他能赦免我的罪，并把那些灵魂赐给我。}\OriginalTerm{主帅}{chief-captain}听了他的话，他们在主面前恳求。他们呼求了很久，有声音从天上说：\Quote{\Person{亚巴郎}，\Person{亚巴郎}，我垂听了你的声音和你的祈祷，赦免了你的罪。你以为我毁灭的那些人，我已因我的极大仁慈把他们唤回并带入生命，因为我在审判中暂时惩罚了他们。而那些我在地上毁灭的活人，我不会用死亡来惩罚他们。}\par

15. 主的声音也对统帅\Person{弥额尔}说：\Quote{\Person{弥额尔}，我的仆人，将\Person{亚巴郎}带回他的家中，因为看，他的结局已近，他的生命之量已满，使他可以安排一切，然后带他来见我。}于是统帅调转战车与云彩，将\Person{亚巴郎}带回他的家中。\Person{亚巴郎}进入他的内室，坐在他的卧榻上。他的妻子\Person{撒辣}前来拥抱无形体者的脚，谦卑地说话：\Quote{我的主，我感谢祢，因祢带回了我的主\Person{亚巴郎}，因为我们以为他已从我们中被提去了。}他的儿子\Person{依撒格}也前来，伏在他的颈上。同样，他所有的男仆与女仆都围绕\Person{亚巴郎}，拥抱他，光荣天主。无形体者对他们说：\Quote{义人\Person{亚巴郎}，请听。看，你的妻子\Person{撒辣}，看，你的爱子\Person{依撒格}，看，你所有的男仆与女仆都围绕着你。安排你所有的一切，因为日子已近，你将要离开肉体，一次而永远地到主那里去。}\Person{亚巴郎}说：\Quote{这是主说的，或是你自己说的？}统帅回答：\Quote{义人\Person{亚巴郎}，请听。主已命令，我告诉你。}\Person{亚巴郎}说：\Quote{我不与你同去。}统帅听了这些话，立刻从\Person{亚巴郎}面前出去，升到天上，站在至高天主面前，说：\Quote{上主全能者，看，我听从了祢的朋友\Person{亚巴郎}向祢所说的一切，并实现了他的请求。我向他展示了祢的大能，以及天下的一切大地与海洋。我藉云彩与战车向他展示了审判与报应，而他又说：\InnerQuote{我不与你同去。}}至高者对天使说：\Quote{我的朋友\Person{亚巴郎}竟又说：\InnerQuote{我不与你同去？}}天使长说：\Quote{上主全能者，他这样说，我克制自己不动手抓他，因为他从起初就是祢的朋友，且行事都合祢的眼目。地上没有像他的人，连奇人\Person{约伯}也没有，因此我克制自己不动手抓他。永生之王，请命令，该做什么。}\par

16. 于是至高者说：\Quote{给我召那被称为无耻面貌与无情眼目的死亡来。}无形体者\Person{弥额尔}就去对死亡说：\Quote{前来，创造之主、永生之王召你。}死亡听了，战栗发抖，被极大的恐惧所占据，带着极大的畏惧前来，站在无形之父面前，颤抖、呻吟、战栗，等候主的命令。于是无形天主对死亡说：\Quote{前来，你这世界苦涩而凶残的名字，隐藏你的凶残，掩盖你的腐朽，从你身上抛去你的苦涩，穿上你的美丽与一切荣耀，下到我朋友\Person{亚巴郎}那里，带他来见我。但如今我也告诉你，不要惊吓他，只用温和的话语带他来，因为他是我自己的朋友。}死亡听了，从至高者面前出去，穿上极其光明的长袍，使他的容貌如同太阳，变得比世人更美更俊，取了天使长的形像，双颊如火焰燃烧，他出发到\Person{亚巴郎}那里。义人\Person{亚巴郎}从内室出来，坐在\Place{玛默勒}的橡树下，手托下巴，等候天使长\Person{弥额尔}的来临。看，一阵甘甜的香气与一道闪光来到他那里，\Person{亚巴郎}转身看见死亡带着极大的荣耀与美丽向他走来。\Person{亚巴郎}起身迎上前去，以为那是天主的统帅。死亡看见他，向他致意：\Quote{珍贵的\Person{亚巴郎}，义人的灵魂，至高天主的真朋友，圣天使的同伴，愿你喜乐。}\Person{亚巴郎}对死亡说：\Quote{愿你喜乐，容貌与形如同太阳、最荣耀的帮助者、光明带来者、奇妙的人，你的荣耀从何来到我们这里？你是谁？从何而来？}死亡说：\Quote{最义的\Person{亚巴郎}，看，我告诉你真理。我是死亡之苦。}\Person{亚巴郎}对他说：\Quote{不，你是世界的俊美，你是天使与人的荣耀与美丽，你的形比一切更美，竟说你是死亡之苦，而不说你比一切美物更美？}死亡说：\Quote{我告诉你真理。主如何称呼我，我就照样告诉你。}\Person{亚巴郎}说：\Quote{你为何来此？}死亡说：\Quote{我是为你的圣灵魂而来。}\Person{亚巴郎}说：\Quote{我知道你的意思，但我不与你同去。}死亡沉默，一言不答。\par

17. 那时\Person{亚伯拉罕}起身，进了他的房屋，死亡也跟随他进去。\Person{亚伯拉罕}上到他的内室，死亡也与他同上去。\Person{亚伯拉罕}躺在卧榻上，死亡前来坐在他的脚旁。\Person{亚伯拉罕}说：\Quote{离开，离开我，因为我想要在卧榻上歇息。}死亡说：\Quote{我不离开，直到我从你那里夺走你的灵魂。}\Person{亚伯拉罕}对他说：\Quote{我指着永生天主命令你，告诉我真相。你是死亡吗？}死亡对他说：\Quote{我是死亡。我是世界的毁灭者。}\Person{亚伯拉罕}说：\Quote{我恳求你，既然你是死亡，请告诉我，你来到众人面前是否都这样公平、光荣而美丽？}死亡说：\Quote{不，我主\Person{亚伯拉罕}，因为你的义德、你无边的款待之海、以及你对天主的伟大爱情，已成为我头上的冠冕，我在美丽、极大的平安与温和中接近义人；但对于罪人，我带着极大的腐朽、凶猛、最大的痛苦，以及凶猛而无情的神色而来。}\Person{亚伯拉罕}说：\Quote{我恳求你，听我说，向我显示你的凶猛和你一切的腐朽与痛苦。}死亡说：\Quote{你无法看见我的凶猛，至义的\Person{亚伯拉罕}。}\Person{亚伯拉罕}说：\Quote{是的，我必能藉着生活的天主之名看见你一切的凶猛，因为我在天之天主的威能与我同在。}于是死亡脱下他所有的俊美与美丽，他的一切光荣，以及他像太阳般的形像，给自己穿上暴君的长袍，使他显得阴郁，比一切野兽更凶猛，比一切不洁更不洁。他向\Person{亚伯拉罕}显示七条火蛇的头和十四个面孔：一个是火焰与极凶猛的，一个是黑暗的面孔，一个是最阴郁的毒蛇面孔，一个是最可怕的悬崖面孔，一个比蝰蛇更凶猛的面孔，一个可怕的狮子面孔，一个角蝰与蛇怪的面孔。他又给他显示一个火弯刀的面孔，一个持剑的面孔，一个闪电的面孔，可怕地闪着电光，以及恐怖的雷声。他又给他显示另一个凶猛风暴海的面孔，一条凶猛急流的江河，一条可怕的三头蛇，以及一个混有毒药的杯；简言之，他给他显示极大的凶猛和难以忍受的痛苦，以及每一种带死亡气味的必死疾病。由于这极大的痛苦与凶猛，约有七千男仆和女仆死去，而义人\Person{亚伯拉罕}陷入对死亡的麻木，以致他的灵魂无力。\par

18. 至圣的\Person{亚伯拉罕}看见这些事，就对死亡说：\Quote{我恳求你，全毁灭的死亡，隐藏你的凶猛，穿上你先前拥有的美丽和形状。}死亡立刻隐藏了他的凶猛，穿上了他先前的美丽。\Person{亚伯拉罕}对死亡说：\Quote{你为什么这样做，以致杀死了我所有的男女仆役？天主今天为此目的派遣你到这里来吗？}死亡说：\Quote{不，我主\Person{亚伯拉罕}，并非如你所说，而是因你的缘故我被派到这里。}\Person{亚伯拉罕}对死亡说：\Quote{那么这些人是怎么死的？主岂没有吩咐吗？}死亡说：\Quote{请相信，至义的\Person{亚伯拉罕}，这也是一件奇事，你也没有与他们一同被带走。尽管如此，我告诉你真相：因为那时若天主的右手没有与你同在，你也必得离开今生。}义人\Person{亚伯拉罕}说：\Quote{现在我知道我已陷入对死亡的麻木，以致我的灵魂无力；但我恳求你，全毁灭的死亡，既然我的仆役未及寿数而死，来，让我们一起向主、我们的天主祈祷，愿他垂听我们，使那些因你的凶猛未及寿数而死的人复活。}死亡说：\Quote{阿们，但愿如此。}于是\Person{亚伯拉罕}起身，俯伏在地祈祷，死亡也与他一同祈祷；主就派遣生命之神降到死者身上，他们就复活了。于是义人\Person{亚伯拉罕}光荣了天主。\par

19. 亚巴郎上到他的房间躺下，\OriginalTerm{死亡}{Death}前来站在他面前。亚巴郎对他说：\Quote{离开我，因为我想要休息，因为我的精神处于冷漠之中。}死亡说：\Quote{我不会离开你，直到我取走你的灵魂。}亚巴郎面带严峻、怒视着死亡说：\Quote{谁命令你说这话？你自高自大地说了这些话，我不跟你去，直到队长\Person{弥额尔}来到我这里，我才跟他去。但我也告诉你，如果你想要我跟你同行，就向我解释你所有的变化：那七个蛇的火头、悬崖的面貌、锋利的剑、咆哮的大河、狂暴翻腾的大海。也教导我那不可忍受的雷声、可怖的闪电、以及掺了毒药的恶臭杯。教导我这一切。}死亡回答说：\Quote{正义的亚巴郎，请听。七个时代以来，我毁灭世界，将所有人——君王与统治者、富人与穷人、奴隶与自由人——都带下\OriginalTerm{阴府}{Hades}，我护送他们到阴府深处，因此我给你看七个蛇头。我给你看火的面貌，因为许多人被火烧死，看吧，死亡通过火的面貌。我给你看悬崖的面貌，因为许多人从树顶或可怕的悬崖上坠落，失去生命，看到死亡如同可怕的悬崖。我给你看剑的面貌，因为许多人在战争中被剑杀死，看到死亡如同剑。我给你看大河奔流的面貌，因为许多人被大水卷入，被大河冲走，淹死丧命，在未到之时见到死亡。我给你看愤怒翻腾的大海的面貌，因为许多人在海中遭遇巨浪，船毁人亡，被吞噬，看到死亡如同大海。我给你看不可忍受的雷声和可怖的闪电，因为许多人在愤怒的时刻遭遇不可忍受的雷声和可怖的闪电，被抓住，从而见到死亡。我也给你看毒兽：毒蛇、蛇怪、豹子、狮子、幼狮、熊、毒蛇，总之，最正义者，我给你看每一种野兽的面貌，因为许多人被野兽毁灭，另一些人被毒蛇、蛇、毒蛇、角蝰、蛇怪、毒蛇咬伤，断气而死。我也给你看掺了毒药的毁灭之杯，因为许多人被他人下毒，立刻出乎意料地离世。}\par

20. 亚巴郎说：\Quote{我恳求你，还有意外的死亡吗？请告诉我。}死亡说：\Quote{我实实在在地在天主的真理中告诉你，有七十二种死亡。其中一种是正义的死亡，在指定的时刻来临；许多人在一个时辰内进入死亡，被埋葬。看，我已告诉你你所问的一切，现在我对你说，最正义的亚巴郎，放弃一切计谋，不要再问任何事情，一次了结，来，跟我走，如同天主和审判万物的那一位所命令我的。}亚巴郎对死亡说：\Quote{再离开我一会儿，让我在床上休息，因为我心里非常衰弱；自从我亲眼见到你，我的力气耗尽，我肉体的所有肢体对我来说如同铅块，我的精神极其痛苦。离开一会儿，因为我说过我无法忍受看见你的形状。}于是他的儿子\Person{依撒格}前来，伏在他胸前哭泣；他的妻子\Person{撒辣}前来拥抱他的脚，悲痛地哀哭。他的男仆和女仆也前来围住他的床，放声痛哭。亚巴郎进入了死亡的冷漠状态，死亡对亚巴郎说：\Quote{来，握住我的右手，愿喜乐、生命和力量临到你。}因为死亡欺骗了亚巴郎，他握住了死亡的右手，立刻他的灵魂粘附在死亡的手上。与此同时，总领天使\Person{弥额尔}带着一大群天使前来，用神圣编织的细麻布接住他宝贵的灵魂，他们用神圣的香膏和香料照料正义的亚巴郎的身体，直到他死后第三天，将他埋葬在应许之地，\OriginalTerm{玛默勒}{Mamre}的橡树那里；天使们则接住他宝贵的灵魂，升到天上，向万有的主、天主咏唱《\Quote{三圣}》颂歌，并将它安置在那里朝拜天主和圣父。在向主献上极大的赞美和荣耀之后，亚巴郎俯伏朝拜，这时天主和圣父那无玷的声音如此说：\Quote{将我的朋友亚巴郎带到\OriginalTerm{乐园}{Paradise}去，那里有我义人的帐幕，以及我的圣人\Person{依撒格}和\Person{雅各伯}在他怀中的住所，那里没有烦恼、没有忧愁、没有叹息，只有平安、喜乐和永生。}（亲爱的弟兄们，让我们也效法圣祖亚巴郎的好客，达到他德行的生活方式，好使我们堪当承受永生，光荣圣父、圣子及圣神；愿光荣与权能归于他，直到永远。阿们。）\par

% structure: p0023 source=h2 target=section reason=relative-heading-rank; base=h2

\section{版本二}

1. 当亚巴郎死期临近时，主对\Person{弥额尔}说：\Quote{起来，到我的仆人亚巴郎那里去，对他说：\InnerQuote{你将离开生命，因为看，你现世生命的时日已满；好使他在死前料理好他的家事。}}\par

2. 弥额尔去了，来到亚巴郎那里，见他正坐在耕牛前犁地，他外貌极其苍老，怀里抱着他的儿子。亚巴郎看见总领天使弥额尔，从地上起身向他行礼，并不认识他是谁，便对他说：\Quote{愿主保佑你。愿你的旅途平安。}弥额尔回答他说：\Quote{你真是善良的好父亲。}亚巴郎回答他说：\Quote{来吧，靠近我，兄弟，请稍坐片刻，我好吩咐人牵一头牲口来，我们好回我家去，你在我这里歇息，因为天色已晚，到早晨你再起来，往你愿意去的地方去，免得有恶兽遇见你伤害你。}弥额尔问亚巴郎说：\Quote{请告诉我你的名字，免得我进你家门后成为你的负担。}亚巴郎回答说：\Quote{我的父母叫我亚巴郎，主却给我起名叫亚巴郎，说：\InnerQuote{起来，离开你的家、你的亲族，往我将要指示你的地方去。}当我前往主所指示我的地方时，他对我说：\InnerQuote{你的名字不再叫亚巴郎，而要叫亚巴郎。}}弥额尔回答他说：\Quote{请宽恕我，我的父，有经验的天主的人，因为我是外乡人，我听说你曾走了四十斯塔狄的路，牵了一只山羊宰了，在你家中款待了天使，使他们得以在那里休息。}他们这样说着话，便起身向屋子走去。亚巴郎叫了一个仆人对他说：\Quote{去，给我牵一头牲口来，让这位外乡人骑上，因为他旅途劳累了。}弥额尔说：\Quote{不要麻烦那年轻人，我们轻松地走，直到我们到家，因为我喜爱与你为伴。}\par

3. 他们起身继续前行，快接近城时，离城大约三斯塔狄，他们看见一棵大树，有三百根枝子，像柽柳树。他们听见从枝子间传来歌声：\Quote{你是圣的，因为你持守了你被派遣的目的。}亚巴郎听见这声音，将奥秘藏在心中，心里说：\Quote{我所听见的这是什么奥秘呢？}他进了屋子，对仆人说：\Quote{起来，到羊群中去，牵三只羊来，快快宰了，预备好，让我们吃喝，因为今天是我们的喜庆日。}仆人们牵了羊来，亚巴郎叫他的儿子依撒格说：\Quote{我儿依撒格，起来拿水到盆里，让我们给这位外乡人洗脚。}依撒格照吩咐做了，亚巴郎说：\Quote{我看，从此以后，我再也不会用这个盆给任何来到我们这里的客人洗脚了。}依撒格听见父亲这样说，就哭着对他说：\Quote{父亲，你这是说什么呢？这是我给外乡人洗脚的最后一次吗？}亚巴郎见儿子哭泣，也大大地哭泣，弥额尔见他们哭泣，也哭了，弥额尔的眼泪掉在盆里，变成了一块宝石。\par

4. 撒辣在屋里听见他们的哭声，就出来对亚巴郎说：\Quote{主啊，你们为什么这样哭呢？}亚巴郎回答她说：\Quote{没有坏事。你回屋去，做你自己的事，免得我们打扰这个人。}撒辣便回去了，预备晚饭。太阳快要落山了，弥额尔出了屋子，被提到天上，在天主面前朝拜，因为日落时所有天使都朝拜天主，弥额尔本人是天使中的首位。他们都朝拜了他，各自回到自己的位置。弥额尔却在主面前发言说：\Quote{主啊，请命我在你的圣荣前受问！}主对弥额尔说：\Quote{你宣布什么都可以！}总领天使回答说：\Quote{主啊，你曾派遣我到亚巴郎那里，对他说：\InnerQuote{离开你的身体，离开这世界；主召唤你。}但我不敢，主啊，向他显示我自己，因为他是你的朋友，是义人，是款待客旅的人。但我恳求你，主啊，请命令亚巴郎死亡的记忆进入他自己的心中，而不要命令我告诉他，因为说\InnerQuote{离开世界}，尤其是离开自己的身体，是太突然了，因为你从起初造他，就是为了怜恤所有人的灵魂。}于是主对弥额尔说：\Quote{起来，到亚巴郎那里去，与他同住，你见他吃什么，也吃什么，他在哪里睡，你也在哪里睡。因为我要将亚巴郎死亡的意念藉着梦注入他儿子依撒格的心中。}\par

5. 于是那天晚上弥额尔进了亚巴郎的家，见他们正在预备晚饭，他们吃喝作乐。亚巴郎对他儿子依撒格说：\Quote{起来，我儿，为这个人铺床，让他睡觉，把灯放在灯台上。}依撒格照父亲的吩咐做了，依撒格对父亲说：\Quote{我也要来睡在你旁边。}亚巴郎回答他说：\Quote{不，我儿，免得我们打扰这个人，你回自己的房间去睡吧。}依撒格不愿违抗父亲的命令，就去睡在自己的房间里了。\par

6. 约在夜间的第七时辰，\Person{依撒格}醒来，来到父亲寝室的门口，哭着喊道：\Quote{父亲，开门，让我触摸你，趁他们把你从我这里带走之前。} \Person{亚巴郎}起身为他开了门，\Person{依撒格}进去，抱住父亲的脖子哭泣，并哀伤地亲吻他。\Person{亚巴郎}与他的儿子一同哭泣，\Person{弥额尔}看见他们哭泣，也跟着哭了。\Person{撒辣}听见他们哭泣，从卧室里喊道：\Quote{我的主\Person{亚巴郎}，为什么这样哭泣？那位陌生人是否告诉了你关于你兄弟的儿子\Person{罗特}的死讯？还是有什么事临到了我们？} \Person{弥额尔}回答\Person{撒辣}说：\Quote{不，\Person{撒辣}，我没有带来\Person{罗特}的消息，但我知晓你的一切善心，在这方面你超越世上所有的人，主已记念了你。} 于是\Person{撒辣}对\Person{亚巴郎}说：\Quote{当天主的人来到你这里时，你怎么敢哭泣？你的眼睛为何流泪，因为今日有极大的喜乐。} \Person{亚巴郎}对她说：\Quote{你怎么知道这是天主的人？} \Person{撒辣}回答说：\Quote{因为我可以说并宣告，这就是曾在我们家住于\Place{玛默勒}橡树旁的那三个人之一；那时一个仆人前去取了一只小山羊，你宰杀了它，并对我说：\InnerQuote{起来，准备，好让我们与这些人在家中吃饭。} } \Person{亚巴郎}回答说：\Quote{你观察得很对，妇人！因为当我洗他的脚时，心里也知道这就是我在\Place{玛默勒}橡树旁所洗的脚；当我开始询问他的行程时，他对我说：\InnerQuote{我去保全你的兄弟\Person{罗特}，脱离\Place{索多玛}人。}于是我便明白了这奥秘。}\par

7. \Person{亚巴郎}对\Person{弥额尔}说：\Quote{天主的人，请告诉我，并让我知道你为什么来到此地。} \Person{弥额尔}说：\Quote{你的儿子\Person{依撒格}会告诉你。} \Person{亚巴郎}对他的儿子说：\Quote{我亲爱的儿子，告诉我你今天在梦中看见了什么，为什么你受惊了？讲给我听。} \Person{依撒格}回答他的父亲说：\Quote{我在梦中看见了太阳和月亮，我头上有一顶冠冕；从天上来了一个身材高大的人，照耀如光，被称为光之父。他从我头上取走了太阳，却将光线留给了我。我哭着说：\InnerQuote{我恳求你，我主，不要取走我头上的荣耀、我家的光明和我的一切荣耀。}太阳、月亮和星辰也都哀叹说：\InnerQuote{不要取走我们力量的荣耀。}那发光的人回答我说：\InnerQuote{不要因我取走你家的光明而哭泣，因为它是从患难中被提到安息，从卑微升到崇高；他们把他从狭窄处举到宽阔处，从黑暗提升到光明。}我对他说：\InnerQuote{我恳求你，主，也把光线一同取去。}他对我说：\InnerQuote{白天有十二个时辰，那时我将取走所有的光线。}发光的人说这话时，我看见我家的太阳升上天去，但那冠冕我再也没有看见；那太阳就像你，我的父亲。} \Person{弥额尔}对\Person{亚巴郎}说：\Quote{你的儿子\Person{依撒格}说的是真话，因为你将要去，并被提升到天上，但你的身体要留在地上，直到七千个时代期满，那时所有血肉之躯都将复活。所以现在，\Person{亚巴郎}，安排你的家和你的子女，因为你已完全听到了关于你所定的旨意。} \Person{亚巴郎}回答\Person{弥额尔}说：\Quote{我恳求你，主，如果我必须离开我的身体，我渴望带着我的身体被提升，好让我看到我主天主在天上和地上所造的万物。} \Person{弥额尔}回答说：\Quote{这事不是我所能做的，但我要去告诉主这事，如果蒙他吩咐，我就把这一切都指示给你。}\par

8. \Person{弥额尔}升到天上，在主面前论及\Person{亚巴郎}。主回答\Person{弥额尔}说：\Quote{去，带着\Person{亚巴郎}的身体把他提上去，指示他一切事物；无论他对你说什么，你都要为他做，如同对待我的朋友。}于是\Person{弥额尔}出去，在云上带着\Person{亚巴郎}的身体把他提上去，带他到了大洋河。\par

12. \Person{亚巴郎}看见了审判之处后，云彩把他降到下面的穹苍上。\Person{亚巴郎}俯视大地，看见一个人与有夫之妇通奸。\Person{亚巴郎}转身对\Person{弥额尔}说：\Quote{你看见这邪恶吗？主，请从天上降下火来烧灭他们。}立刻有火降下烧灭了他们，因为主曾对\Person{弥额尔}说：\Quote{凡\Person{亚巴郎}求你为他做的，你都要做。} \Person{亚巴郎}再观看，看见另一些人辱骂他们的同伴，就说：\Quote{让地裂开吞没他们。}他一说话，地就裂开活活吞没了他们。云彩又领他到另一个地方，\Person{亚巴郎}看见一些人进入旷野去杀人，就对\Person{弥额尔}说：\Quote{你看见这邪恶吗？让野兽从旷野出来，把他们撕碎。}同一时辰，野兽从旷野出来，吞噬了他们。于是主天主对\Person{弥额尔}说：\Quote{带\Person{亚巴郎}回他自己的家，不要让他周游我所造的全部受造物，因为他对罪人没有怜悯心，但我怜悯罪人，好使他们悔改而存活，并因他们的罪悔改而得救。}\par

8. \Person{亚巴郎}观看，看见两扇门，一扇小，一扇大。两门之间坐着一个人，坐在一个极其荣耀的宝座上，周围有许多天使；他时而哭泣，时而欢笑，但他的哭泣超过他的欢笑七倍。\Person{亚巴郎}对\Person{弥额尔}说：\Quote{这位是谁？他坐在两门之间，极其荣耀；他有时欢笑，有时哭泣，他的哭泣超过他的欢笑七倍。} \Person{弥额尔}对\Person{亚巴郎}说：\Quote{你不知道他是谁吗？}他说：\Quote{主，我不知道。} \Person{弥额尔}对\Person{亚巴郎}说：\Quote{你看见这两扇门，一扇小，一扇大吗？它们引导人走向生命和毁灭。坐在它们之间的这个人是\Person{亚当}，主所创造的第一个人，并把他安置在此处，让他观看每个离开身体的灵魂，因为众人都是从他而出。因此，当你看见他哭泣时，要知道他看见许多灵魂被带向毁灭；但当你看见他欢笑时，他看见许多灵魂被带向生命。你看见他的哭泣如何超过他的欢笑吗？因为他看见世上大多数的人通过宽门被带向毁灭，所以他的哭泣超过他的欢笑七倍。}\par

9. \Person{亚巴郎}说：\Quote{那不能从窄门进入的人，就不能进入生命吗？}于是\Person{亚巴郎}哭着说：\Quote{我可悲啊！我该怎么办？因为我是一个身体宽阔的人，怎能进入那连十五岁少年都无法进入的窄门呢？}\Person{弥额尔}回答\Person{亚巴郎}说：\Quote{父亲，不要害怕，也不要悲伤，因为你必能毫无阻碍地进入，所有像你一样的人也是如此。}当\Person{亚巴郎}站着惊讶时，看，有一位主的天使正驱赶六万个罪人的灵魂去灭亡。\Person{亚巴郎}对\Person{弥额尔}说：\Quote{所有这些人都要进入灭亡吗？}\Person{弥额尔}对他说：\Quote{是的，但我们去这些灵魂中查寻一下，看看其中是否有一个义人。}他们去了，发现一位天使手里拿着从这六万个灵魂中的一个妇女的灵魂，因为他发现她的罪与她的所有工作重量相等，既不活动也不静止，而是处于\OriginalTerm{中间状态}{state between}；但其余的灵魂被带往灭亡。\Person{亚巴郎}对\Person{弥额尔}说：\Quote{主啊，这就是那将灵魂从身体中带出的天使吗？}\Person{弥额尔}回答说：\Quote{这是死亡，他将他们带入审判的地方，让审判者审判他们。}\par

10. \Person{亚巴郎}说：\Quote{我的主，我恳求您带我到审判的地方，让我也能看到他们如何受审。}于是\Person{弥额尔}将\Person{亚巴郎}带到一片云上，领他进入乐园。当他来到审判者所在的地方时，那位天使前来把那灵魂交给审判者。灵魂说：\Quote{主，求祢怜悯我。}审判者说：\Quote{我怎能怜悯你？当你对你所生的女儿，你腹中的果子毫无怜悯时，你为何杀了她？}灵魂回答说：\Quote{不，主，杀人之事并非我所为，是我的女儿诬陷了我。}但审判者命令那记录案卷的人前来，看，有革鲁宾拿着两本书。与他们一起的有一个身材极其高大的人，头上戴着三顶冠冕，其中一顶冠冕比另外两顶更高。这些称为\OriginalTerm{见证之冠}{crowns of witness}。那人手里拿着一支金笔，审判者对他说：\Quote{显明这灵魂的罪。}那人打开革鲁宾的一本书，查寻那妇女灵魂的罪，就找着了。审判者说：\Quote{悲惨的灵魂，你为何说没有犯杀人罪？在你丈夫死后，你不是去与你女儿的丈夫通奸，并杀了她吗？}他又定了她其他的罪，就是她自幼年所犯的一切罪。那妇女听见这些话，喊道：\Quote{我真可悲！我在世上所犯的一切罪我都忘记了，但在这里却没有被忘记。}于是他们也将她带走，交给施刑者。\par

11. \Person{亚巴郎}对\Person{弥额尔}说：\Quote{主啊，这位审判者是谁？那另一位定罪的又是谁？}\Person{弥额尔}对\Person{亚巴郎}说：\Quote{你看见那审判者吗？他是\Person{亚伯尔}，他是第一个作证的人，\Person{天主}将他带到这里来审判。而在这里作证的是天地之师、正义之书记\Person{哈诺客}，因为主派他们到这里来记录每个人的罪和义。}\Person{亚巴郎}说：\Quote{\Person{哈诺客}未经历过死亡，怎能承担灵魂的重担？他怎能给所有的灵魂判决？}\Person{弥额尔}说：\Quote{他并不判决灵魂，这是不允许的。但\Person{哈诺客}自己并不判决，而是主在判决，他只负责记录。因为\Person{哈诺客}曾向主祈祷说：\InnerQuote{主，我不愿判决灵魂，免得我使任何人感到沉重。}主对\Person{哈诺客}说：\InnerQuote{我命令你记录那做补赎的灵魂的罪，它必进入生命；若那灵魂不做补赎也不悔改，你必发现它的罪被记录下来，它必被投入惩罚。}}约在第九时辰，\Person{弥额尔}将\Person{亚巴郎}带回他的家。但他的妻子\Person{撒辣}，因不知\Person{亚巴郎}去了何处，忧伤过度，断了气。\Person{亚巴郎}回来后发现她已死，便埋葬了她。\par

13. 当\Person{亚巴郎}死亡的日子临近时，主天主对\Person{弥额尔}说：\Quote{\Person{死亡}不敢走近取走我仆人的灵魂，因为他是我的朋友。但你去将\Person{死亡}装扮得极为美丽，这样派他到\Person{亚巴郎}那里，使他能用眼睛看见\Person{死亡}。}\Person{弥额尔}立刻按所吩咐的，将\Person{死亡}装扮得极为美丽，这样派他到\Person{亚巴郎}那里，使他能看见\Person{死亡}。\Person{死亡}坐在\Person{亚巴郎}附近，\Person{亚巴郎}看见\Person{死亡}坐在他旁边，十分害怕。\Person{死亡}对\Person{亚巴郎}说：\Quote{万福，圣洁的灵魂！万福，主天主的朋友！万福，旅客的安慰和款待者！}\Person{亚巴郎}说：\Quote{至高天主的仆人，欢迎你。我恳求你，告诉我你是谁；然后进入我家，享用饮食，然后离开我，因为我从看见你坐在我旁边，我的灵魂就烦乱。因为我根本不配接近你，因你是崇高的神，而我是血肉之躯，所以我不能承受你的荣耀，因我看你的美不属于这个世界。}\Person{死亡}对\Person{亚巴郎}说：\Quote{我告诉你，在天主所造的一切受造物中，没有找到一个像你一样的，因为主亲自查寻，在全地上也找不到这样的一个。}\Person{亚巴郎}对\Person{死亡}说：\Quote{你怎么敢说谎？因为我看你的美不属于这个世界。}\Person{死亡}对\Person{亚巴郎}说：\Quote{亚巴郎，不要以为这美是我的，或以为我如此来到每个人面前。不，若有人像你一样是义人，我就这样戴着冠冕来到他那里；但若是罪人，我就以大腐败而来，并以他们的罪为我做冠冕戴在头上，用极大的恐惧摇动他们，使他们惊慌失措。}\Person{亚巴郎}于是对他说：\Quote{那你的美从何而来？}\Person{死亡}说：\Quote{没有比我更充满腐败的了。}\Person{亚巴郎}对他说：\Quote{你真是那被称为死亡的吗？}他回答说：\Quote{我就是那\OriginalTerm{苦涩之名}{bitter name}。我是哭泣……}\par

14. 亚巴郎对死亡说：\Quote{请向我们显示你的腐朽。}死亡便显露出他的腐朽；他有两个头，一个有如毒蛇的面孔，凡被毒蛇咬死的人即因它而死；另一个头像一把剑，凡死于刀剑或弓箭的人即因它而死。那一天，亚巴郎的仆人因见死亡而恐惧而死。亚巴郎见状，便向主祈祷，主使他们复活。天主又返回，如同在梦中取去亚巴郎的灵魂，总领天使弥额尔便将它带到天上。依撒格将他葬在母亲撒辣旁边，光荣赞美天主，因为光荣、尊威和朝拜归于他，圣父、圣子及圣神，从今世直到永远。阿们。\par

\SourceRecord{Translated by W.A. Craigie.；Ante-Nicene Fathers；Vol. 9.；Edited by Allan Menzies.；Buffalo, NY: Christian Literature Publishing Co.,；1896.；<http://www.newadvent.org/fathers/1007.htm>.}
